\section{Continuous Assessment 2 (CA-2) --- Set A}
\label{sec:ca2_set_a}

\LPUPaperHeader{Continuous Assessment (CA-2)}{A}{60 Minutes}{30 Marks}{CO4 (NumPy Arrays), CO5 (Pandas \& Matplotlib)}

\begin{center}
\sffamily\textbf{\color{AlertRed}Note: Attempt any five out of the six questions. Each question carries 6 marks.}
\end{center}

\questionitem{1}{
    Differentiate between a Python \texttt{list} and a NumPy \texttt{ndarray} in terms of memory layout, type heterogeneity, and SIMD hardware acceleration.
}{6}{CO4}{L2 Understand}

\questionitem{2}{
    Explain the \textbf{Three Universal Rules of NumPy Broadcasting}. Determine whether array shapes $(8, 1, 6, 1)$ and $(7, 1, 5)$ are broadcast-compatible, and calculate the resulting shape.
}{6}{CO4}{L3 Apply}

\questionitem{3}{
    Write a Python function using NumPy to compute the \textbf{moving average} of a 1D signal with window size $k$ using \texttt{np.convolve()}. Compare its asymptotic time complexity against a naive nested loop.
}{6}{CO4}{L4 Analyze}

\questionitem{4}{
    Differentiate between label-based indexing (\texttt{.loc}) and integer-position indexing (\texttt{.iloc}) in Pandas. Explain the critical difference in slicing endpoint inclusion.
}{6}{CO5}{L2 Understand}

\questionitem{5}{
    Write a Pandas pipeline that reads an employee dataset, handles missing values by imputing departmental median salaries, and creates a pivot table showing total payroll by department and employment tier.
}{6}{CO5}{L3 Apply}

\questionitem{6}{
    Using Matplotlib's object-oriented interface (\texttt{plt.subplots}), write a script to produce a 2-panel figure showing a line plot with error bars on the left and a normalized probability histogram on the right.
}{6}{CO5}{L3 Apply}

\newpage
\subsection*{Solutions and Marking Rubrics --- CA-2 Set A}

\begin{solutionbox}[Solution to Question 1]
\begin{itemize}[leftmargin=1.5em]
    \item \textbf{Memory Architecture}: Python lists store arrays of 64-bit pointers to disjoint heap objects (\texttt{PyObject*}), requiring pointer dereferencing and causing CPU cache misses. NumPy \texttt{ndarray} stores a contiguous block of homogeneous binary data with fixed stride offsets.
    \item \textbf{Type Enforcement}: Lists allow arbitrary mixed types; \texttt{ndarray} requires a single homogeneous \texttt{dtype}.
    \item \textbf{Vectorization}: NumPy operations execute via compiled C/Fortran libraries leveraging CPU SIMD (Single Instruction, Multiple Data) registers, yielding 50x--100x speedups.
\end{itemize}
\begin{markingrubric}
$\bullet$ Pointer array vs contiguous binary buffer: \textbf{2 Marks}\\
$\bullet$ Homogeneous dtype requirement: \textbf{2 Marks}\\
$\bullet$ SIMD hardware acceleration explanation: \textbf{2 Marks}
\end{markingrubric}
\end{solutionbox}

\begin{solutionbox}[Solution to Question 2]
\textbf{The Three Broadcasting Rules}:
\begin{enumerate}[leftmargin=1.5em]
    \item Prepend 1s to the shape of the array with fewer dimensions until rank matches.
    \item Compare corresponding dimensions from right to left: they must either be equal, or one of them must be 1.
    \item Arrays with dimension size 1 are stretched along that axis to match the other array.
\end{enumerate}
\textbf{Evaluation of $(8, 1, 6, 1)$ and $(7, 1, 5)$}:
\begin{align*}
\text{Array 1:} \quad & 8 & 1 & 6 & 1 \\
\text{Array 2:} \quad & 1 & 7 & 1 & 5 \quad \text{(prepended with 1)} \\
\midrule
\text{Result:} \quad & \mathbf{8} & \mathbf{7} & \mathbf{6} & \mathbf{5}
\end{align*}
All dimensions are compatible ($8 \times 1 \to 8$, $1 \times 7 \to 7$, $6 \times 1 \to 6$, $1 \times 5 \to 5$). The resulting array shape is $\mathbf{(8, 7, 6, 5)}$.
\begin{markingrubric}
$\bullet$ Formulation of the 3 broadcasting rules: \textbf{3 Marks}\\
$\bullet$ Right-to-left trailing dimension alignment: \textbf{2 Marks}\\
$\bullet$ Final shape calculation $(8, 7, 6, 5)$: \textbf{1 Mark}
\end{markingrubric}
\end{solutionbox}

\begin{solutionbox}[Solution to Question 3]
\begin{pycode}
import numpy as np\\
~\\
def moving\_average(signal, k):\\
\_\_ESCTOKEN\_7\_\_ if k <= 0 or k > len(signal):\\
\_\_ESCTOKEN\_7\_\_ \_\_ESCTOKEN\_7\_\_ raise ValueError("Window k must satisfy 1 <= k <= len(signal)")\\
\_\_ESCTOKEN\_7\_\_ kernel = np.ones(k) / k\\
\_\_ESCTOKEN\_7\_\_ return np.convolve(signal, kernel, mode='valid')
\end{pycode}
\textbf{Complexity}: Naive nested Python loops take $\mathcal{O}(N \cdot k)$ time with dynamic interpretation overhead. NumPy's compiled FFT/C convolution executes in $\mathcal{O}(N)$ or $\mathcal{O}(N \log k)$ time with cache spatial locality.
\begin{markingrubric}
$\bullet$ Implementation via \texttt{np.convolve} with normalized kernel: \textbf{3 Marks}\\
$\bullet$ Parameter validation and \texttt{mode='valid'} boundary: \textbf{1 Mark}\\
$\bullet$ Asymptotic complexity comparison: \textbf{2 Marks}
\end{markingrubric}
\end{solutionbox}

\begin{solutionbox}[Solution to Question 4]
\begin{itemize}[leftmargin=1.5em]
    \item \texttt{.loc}: Label-based selector. Accepts string index labels, date timestamps, or boolean Series. \textbf{Endpoint is INCLUSIVE}: \texttt{df.loc['Jan':'Mar']} includes March.
    \item \texttt{.iloc}: Integer position selector ($0 \dots N-1$). \textbf{Endpoint is EXCLUSIVE}: \texttt{df.iloc[0:3]} includes rows 0, 1, 2 only.
\end{itemize}
\begin{markingrubric}
$\bullet$ Label vs integer coordinate distinction: \textbf{3 Marks}\\
$\bullet$ Inclusive vs exclusive slicing endpoint explanation: \textbf{3 Marks}
\end{markingrubric}
\end{solutionbox}

\begin{solutionbox}[Solution to Question 5]
\begin{pycode}
import pandas as pd\\
~\\
def process\_payroll(filepath):\\
\_\_ESCTOKEN\_7\_\_ df = pd.read\_csv(filepath)\\
\_\_ESCTOKEN\_7\_\_ \_\_ESCTOKEN\_1\_\_ Impute missing salaries with departmental median\\
\_\_ESCTOKEN\_7\_\_ df['salary'] = df.groupby('department')['salary'].transform(\\
\_\_ESCTOKEN\_7\_\_ \_\_ESCTOKEN\_7\_\_ lambda s: s.fillna(s.median())\\
\_\_ESCTOKEN\_7\_\_ )\\
\_\_ESCTOKEN\_7\_\_ \_\_ESCTOKEN\_1\_\_ Pivot table for aggregated payroll\\
\_\_ESCTOKEN\_7\_\_ pivot = pd.pivot\_table(\\
\_\_ESCTOKEN\_7\_\_ \_\_ESCTOKEN\_7\_\_ df, values='salary', index='department',\\
\_\_ESCTOKEN\_7\_\_ \_\_ESCTOKEN\_7\_\_ columns='tier', aggfunc='sum', fill\_value=0\\
\_\_ESCTOKEN\_7\_\_ )\\
\_\_ESCTOKEN\_7\_\_ return pivot
\end{pycode}
\begin{markingrubric}
$\bullet$ Correct use of \texttt{.groupby().transform()} for median imputation: \textbf{3 Marks}\\
$\bullet$ Parameterization of \texttt{pd.pivot\_table}: \textbf{2 Marks}\\
$\bullet$ Handling missing values defensively: \textbf{1 Mark}
\end{markingrubric}
\end{solutionbox}

\begin{solutionbox}[Solution to Question 6]
\begin{pycode}
import matplotlib.pyplot as plt\\
import numpy as np\\
~\\
fig, (ax1, ax2) = plt.subplots(1, 2, figsize=(12, 5))\\
~\\
\# Left panel: Line plot with error bars\\
x = np.linspace(0, 10, 20)\\
y = np.sin(x)\\
y\_err = 0.15 * np.ones\_like(x)\\
ax1.errorbar(x, y, yerr=y\_err, fmt='-o', color='navy', capsize=3)\\
ax1.set\_title("Signal Transmission")\\
ax1.grid(True)\\
~\\
\# Right panel: Normalized histogram\\
samples = np.random.normal(0, 1, 1000)\\
ax2.hist(samples, bins=25, density=True, color='skyblue', edgecolor='black')\\
ax2.set\_title("Density Distribution")\\
~\\
plt.tight\_\_ESCTOKEN\_0\_\_layout()\\
plt.savefig("panel\_output.png", dpi=300)
\end{pycode}
\begin{markingrubric}
$\bullet$ Object-oriented figure and axes instantiation: \textbf{2 Marks}\\
$\bullet$ Correct errorbar plotting syntax: \textbf{2 Marks}\\
$\bullet$ Histogram density normalization and layout export: \textbf{2 Marks}
\end{markingrubric}
\end{solutionbox}

\newpage
